\scanpage{201}
\uline{Proof.} \deleted{Let us fix a structure
\deleted{\unsure{conformal}} \inserted{\deleted{\illegible{1}}} \inserted{complex,}
on $X$, \unsure{invariant} under}

Let us prove d). \deleted{Let $\Gamma_0 =$} We identify
$\Gamma',\Gamma''$ with subgroups of $\Gamma$, and put $\Gamma_0=\Gamma'\cap\Gamma''$.
Let \unsure{$\mathcal E^!$} be the subgroup of $\mathcal E$ consisting of the $g\in\mathcal E$
such that $\operatorname{int}(g).\Gamma_0$ is $\pi$-conjugate to $\Gamma_0$
and whose image in $\Gamma$ \deleted{\illegible{1}} centralizes $\Gamma_0$.
\deleted{\illegible{1}} We have $\mathcal E^!\supset\pi$, and $\mathcal E^!$ is therefore the inverse
image of a subgroup $\Gamma^!$ of $\Gamma$, which
centralizes $\Gamma_0$ \deleted{\illegible{1}}, and which contains $\Gamma'$ and $\Gamma''$
(for $\mathcal E^!$ contains the centralizer of $\Gamma_0$ in $\mathcal E$,
hence $\Gamma'$ and $\Gamma''$). Replacing
$\Gamma$ by $\Gamma^!$, $\mathcal E$ by $\mathcal E^!$ if necessary, we may suppose $\mathcal E=\mathcal E^!$, $\Gamma=\Gamma^!$ i.e.
that $\Gamma_0\subset\mathrm{Centre}$ of $\Gamma$, and that $\Gamma_0\hookrightarrow\mathcal E$
\unsure{invariant} modulo $\pi$-conjugation under $\Gamma$.

\deleted{After this let us choose a topological realization
$(X,S,a)$ of the \deleted{outer} operation of $\Gamma_0$ on $\pi$,
where $a\in U=X-S$ -- thus $\pi\simeq\pi_1(U,a)$. By
the result (\unsure{b}) on existence and uniqueness
of realizations in the cyclic case, this
is possible, and moreover we can have}

\scanpage{202}
\mspage{101}
\deleted{$\Gamma'$ operate on \unsure{$(X,S,a)$} (i.e. on $(U,a)$) in a way
compatible with its action on $\pi_1(U,a)=\pi$,
so as to extend the action of \unsure{$\Gamma_0$|$\Gamma$} --
and \unsure{the same|$=$} for $\Gamma''$. But it is not clear
a priori
that the two actions \deleted{\unsure{can}} are compatible
in the precise sense \deleted{\unsure{of the generated subgroup}}
\inserted{that the \unsure{generated subgroup}}
\deleted{\unsure{by} $\Gamma',\Gamma''$} $\subset\operatorname{Aut\,top}(U)$ is also finite.}

\deleted{Let us now consider $X'=X/\Gamma_0$, $S'=S/\Gamma_0$,
and the image $a'$ of $a$ in $X'$ -- thus $a'\in U'=X'-S'$.
Then \unsure{$\Gamma/\Gamma_0$|$\Gamma'/\Gamma_0$} operates \unsure{isotypically|isotopically} on $(X',S',b)$}

\deleted{Let us use $\pi^{\Gamma_0}=(1)$, we find that}

We are going to construct a section of $\mathcal E$
over $\Gamma$ in \unsure{all} its entirety, as follows. Let $\widetilde\Gamma\subset\mathcal E$
be the centralizer of $\Gamma_0$ in $\mathcal E$, \unsure{we have|$\Rightarrow$}
\deleted{$\gamma\in\Gamma$, $g\in\mathcal E$ over $\gamma$, let us consider}
$\widetilde\Gamma\cap\pi=(1)$ (for this means $\pi^{\Gamma_0}=(1)$)
\deleted{$\operatorname{int}(g).\Gamma_0$, $\exists!\ \alpha\in$}
\inserted{hence}
\deleted{I say that} $\widetilde\Gamma\longrightarrow\Gamma$ injective, I say
that it is surjective. Indeed, let
$\gamma\in\Gamma$, $g\in\mathcal E$ over $\gamma$, by hypothesis
\inserted{on $\Gamma$}
$\exists\ \alpha\in\pi$ such that $\operatorname{int}(g).\Gamma_0=\operatorname{int}(\alpha).\Gamma_0$
i.e. we may suppose $\operatorname{int}(g)\Gamma_0=\Gamma_0$, \deleted{\unsure{hence}} i.e. $g$ normalizes
$\Gamma_0$, but since $\Gamma_0$ is central in $\Gamma$
this means that $g$ centralizes $\Gamma_0$ i.e. $g\in\widetilde\Gamma$.

\scanpage{203}
Thus $\widetilde\Gamma$ \inserted{$\simeq\Gamma$} \unsure{is} a \uline{finite} subgroup of $\operatorname{Aut}(\pi)$,
\deleted{hence cyclic} containing $\Gamma'$ and $\Gamma''$, q.e.d.% a dashed line runs under the deleted words (possible restoration mark); unresolved

\uline{N.B.} If we had not assumed
$\Gamma$ finite at the outset, it would still be true that $\Gamma',\Gamma''$
generate a subgroup $\widetilde\Gamma$ \inserted{$\subset\operatorname{Aut}(\pi)$} such that
$\widetilde\Gamma\cap\pi=(1)$, but that does us a fat lot of good
\deleted{-- \unsure{we can find some},
for example isomorphic to the whole of \unsure{$\operatorname{Autext}(\pi,i)$},
(for a chosen \unsure{$i\in I(\pi)$}: take the
\deleted{\illegible{2}}
centralizer of \unsure{$i$} in \illegible{2} $\operatorname{Aut}(\pi)$)}

\scanpage{204}
\mspage{102}
19. Teichmüller tower.

Let \inserted{$g\in\mathbb N$, and} $X_g$ be a compact orientable surface of
genus $g$, \deleted{\unsure{the set of orientations},}
\unsure{$[$|$]$} $(a_{g,i})_{i\in\mathbb N^*}=(a_{g,1},a_{g,2},\ldots)$ \unsure{$]$|$[$}
a sequence of
distinct points of $X_g$. We put for \unsure{$\forall\nu\in\mathbb N$}
\[
(1)\qquad S_{g,\nu}=\{\unsure{x_{g,i}|a_{g,i}}\mid 0\leq i\leq\nu-1\}
\qquad\text{thus }\nu=\operatorname{card}S_{g,\nu}
\]
\[
(2)\qquad U_{g,\nu}=X_g-S_{g,\nu}.
\]
\margin{N.B. we \unsure{have} $a_{g,\nu}\in U_{g,\nu}$}
We thus have $S_{g,0}=\varnothing$, $U_{g,0}=X_g$; The $S_{g,\nu}$ form
a \inserted{strictly} increasing sequence of finite subsets of $X_g$,
the $U_{g,\nu}$ a strictly decreasing sequence of open subsets
of $X_g$. In what follows we shall take $U_{g,\nu}$ as
the standard surface (oriented, or orientable), of
type $(g,\nu)$.

Let
\[
(3)\qquad A_g=\operatorname{Aut}(X_g)
\]
be the group of automorphisms of $X_g$, equipped with the
topology of uniform convergence of $u$ and of its
inverse. We put
\[
(4)\qquad A_{g,\nu}=\{u\in A_g\mid u(S_{g,\nu})=S_{g,\nu}\}\subset A_g
\]
thus $A_{g,\nu}$ is a subgroup of $A_g$. We
have a canonical isomorphism $u\mapsto u|U_{g,\nu}$
\[
A_{g,\nu}\xrightarrow{\sim}\operatorname{Aut}(U_{g,\nu}).\qquad(5)
\]
\margin{\uline{N.B.} This is \inserted{\unsure{probably|probably}} \unsure{an isomorphism}
topologically, when $A_{g,\nu}$ is
equipped with the topology induced by $A_g$,
and $\operatorname{Aut}(U_{g,\nu})$ with the usual topology of
\unsure{compact} \illegible{1} convergence of $u$ and its inverse.}
\deleted{\unsure{Let}} We have \deleted{\unsure{a}} canonical surjective homomorphism
\[
(6)\qquad A_{g,\nu}\longrightarrow\operatorname{Aut}(S_{g,\nu})
\bigl(\overset{\text{(canonical)}}{\simeq}\unsure{\mathfrak S_\nu}\bigr),
\]

\scanpage{205}
whose kernel is denoted by $A_{g,\nu}^!$
\[
(7)\qquad A_{g,\nu}^!=\{u\in A_g\mid u(a_{g,i})=a_{g,i},\ \forall i\in[0,\ldots,\nu-1]\}
\]
hence exact sequence
\[
(8)\qquad 1\longrightarrow A_{g,\nu}^!\longrightarrow A_{g,\nu}
\longrightarrow\mathfrak S_\nu\longrightarrow1.
\]
Let $A_{g,\nu}^{\circ}$ be the identity component of the group
$A_{g,\nu}$, we thus have
\[
(9)\qquad A_{g,\nu}^{\circ}\bigl(=(A_{g,\nu}^!)^{\circ}\bigr)\subset A_{g,\nu}^!.
\]
Let us put
\[
(10)\qquad \Gamma_{g,\nu}=A_{g,\nu}/A_{g,\nu}^{\circ}=\pi_0(A_{g,\nu})
\quad\text{(Teichmüller group of type }g,\nu\text{)}
\]
\margin{We also put $\Gamma_g=\Gamma_{g,0}(=\Gamma_{g,0}^!)$}
\[
(11)\qquad \Gamma_{g,\nu}^!=A_{g,\nu}^!/A_{g,\nu}^{\circ}=\pi_0(A_{g,\nu}^!).
\]
The exact sequence (8) thus gives an exact
sequence
\[
(12)\qquad 1\longrightarrow\Gamma_{g,\nu}^!\longrightarrow\Gamma_{g,\nu}
\longrightarrow\mathfrak S_\nu\longrightarrow1.
\]
We have canonical homeomorphisms
\[
(13)\qquad A_g/A_{g,\nu}\simeq
\begin{gathered}
\text{open subset }\inserted{\operatorname{Sym}^{\nu}(X)^*}\text{ \inserted{of the} \deleted{\unsure{$X_g$}} symmetric prod-}\\
\text{uct }\inserted{\unsure{\nu^{\text{th}}}}\operatorname{Sym}^{\nu}(X_g)\text{ of }X_g\\
\text{consisting of the finite subsets}\\
\text{of cardinality }\nu\\
=\unsure{\mathcal P_\nu(X_g)}
\end{gathered}
\]
\margin{N.B. the marked points of the \unsure{homogeneous} \unsure{spaces}
on the left-hand sides \unsure{in|of} (13), (14)
\unsure{go to} the systems
of points $\{\unsure{a_{g,0}},a_{g,1},\ldots,a_{g,\nu-1}\}$ or
$(a_{g,0},\ldots,a_{g,\nu-1})$.}
\[
(14)\qquad A_g/A_{g,\nu}^!\simeq
\begin{gathered}
\text{open subset }\inserted{(X_g^{\nu})^*}\text{ of }X_g^{\nu}\text{ consisting of the}\\
\nu\text{-tuples of distinct points}\\
\simeq\operatorname{Mon}(\deleted{\unsure{X_g}}I_\nu,X)\quad\text{(where }I_\nu\\
=\{0,\ldots,\nu-1\}.
\end{gathered}
\]

\scanpage{206}
\mspage{103}
This homeomorphism \unsure{respecting} the natural
actions of $\mathfrak S_\nu\simeq A_{g,\nu}/A_{g,\nu}^!$.

The $A_{g,\nu}^!$ for variable $\nu$ form a
decreasing sequence of groups in $A_g$
\[
(15)\qquad A_g\unsure{=|\supset}A_{g,0}^!\supset A_{g,1}^!\supset A_{g,2}^!
\supset\cdots\supset A_{g,\nu}^!\supset\cdots
\]
and the corresponding morphisms
between homogeneous spaces \inserted{of $A_g$} fit into the
commutative diagram \unsure{of homeomorphisms} \illegible{1}

(16)\quad
\diagram{\begin{tikzcd}[ampersand replacement=\&]
 A_g/A_{g,\nu}^! \arrow[r,"\sim"] \arrow[d] \& \operatorname{Mon}(I_\nu,X) \arrow[d] \\
 A_g/A_{g,\nu'}^! \arrow[r,"\sim"] \& \operatorname{Mon}(I_{\nu'},X)
\end{tikzcd}}
\quad $0\leq\nu'\leq\nu$
\margin{\illegible{1}}
where the second vertical arrow is that of
restriction from $I_\nu$ to $I_{\nu'}\subset I_\nu$.

For $\nu=1$ \inserted{we have $A_{g,1}^!=A_{g,1}$, and} the isomorphism (13) (or if one prefers,
(14)) is written
\[
(17)\qquad A_g/A_{g,1}\simeq X_g
\]
(homeomorphism compatible with \inserted{the} actions of $A_g$)

Moreover if to $\forall x\in X_g$ we associate its
stabilizer $A_{g,x}$ in $A_g$, we find an
obviously surjective map

\scanpage{207}
\[
(18)\qquad X_g\longrightarrow
\begin{gathered}
\text{set of conjugates of the sub-}\\
\text{group }A_{g,1}\text{ of }A_g\\
\simeq A_g/\operatorname{Norm}_{A_g}(A_{g,1})
\end{gathered}
\]
which is identified via (17) with the canonical map
on the homogeneous spaces
\[
 A_g/A_{g,1}\longrightarrow A_g/\operatorname{Norm}_{A_g}(A_{g,1})
\]
deduced from the inclusion $A_{g,1}\subset\operatorname{Norm}_{A_g}(A_{g,1})$,
and we see at once that (18) is
bijective i.e. that \deleted{\illegible{1}}
\[
(19)\qquad A_{g,1}=\operatorname{Norm}_{A_g}(A_{g,1}),
\]
more generally for $\forall\nu$, we have
\[
(20)\qquad A_{g,\nu}=\operatorname{Norm}_{A_g}(A_{g,\nu})
=\operatorname{Norm}_{A_g}(A_{g,\nu}^!)
\]
which means that the canonical maps
of homogeneous $A_g$-sets
\deleted{$x=\{x_1,\ldots,x_\nu\}\longmapsto A_{g,x}=\bigcap$}
\deleted{(23)}
\[
\mathcal P_\nu(X_g)\longrightarrow\text{set of conjugates of }A_{g,\nu}
\]
\[
(21)\qquad S\longmapsto\text{stabilizer }A_{g,S}\text{ of }S
\]
and
\[
(22)\qquad\unsure{\vert}\quad S\longmapsto A_{g,S}^!
\]
are not only surjective, but even bijective.
This comes from the fact that we recover
$S$ in terms of $A_{g,S}$\deleted{${}^!$} \inserted{(or of $A_{g,S}^!$}
\deleted{\unsure{(and} a fortiori in terms \unsure{of the group} $A_{g,S}$)} \inserted{as} \deleted{as}

\scanpage{208}
\mspage{104}
\[
(23)\qquad S=\{x\in X_g\mid ux=x\quad\forall u\in A_{g,S}^!\}
=X_g^{(A_{g,S}^!)}
\]
\[
(24)\qquad\begin{aligned}
S&=\{x\in X_g\mid A_{g,S}.x\text{ finite}\}\\
 &=\{x\in X_g\mid A_{g,S}.x\ne X_g\}
\end{aligned}
\]
\margin{$\forall S\in\mathcal P_\nu(X_g)$}

We also have \unsure{some|\deleted{some}} canonical maps
of homogeneous spaces under $A_g$
\[
(25)\qquad S\longmapsto A_{g,S}^{\circ}
\]
\[
\mathcal P_\nu(X_g)\longrightarrow
\begin{gathered}
\text{set of conjugates}\\
\text{of the }A_{g,\nu}^{\circ}\text{ in }A_g\\
\simeq A_g/\operatorname{Norm}(A_{g,\nu}^{\circ})
\end{gathered}
\]
which is not only surjective, but also
injective, for we have the following relation, which
strengthens \inserted{(23) and} (24), $\forall S\in\mathcal P_\nu(X_g)$
\[
(26)\qquad\begin{aligned}
S&=\{x\in X_g\mid A_{g,S}^{\circ}.x=\{x\}\}=X_g^{A_{g,S}^{\circ}}\\
 &=\{x\in X_g\mid A_{g,S}^{\circ}.x\text{ finite}\}\\
 &=\{x\in X_g\mid A_{g,S}^{\circ}.x\ne X_g\}.
\end{aligned}
\]
\deleted{\illegible{2}} Thus we have for $\forall\nu$
\[
(27)\qquad A_{g,\nu}=\operatorname{Norm}_{A_g}(\unsure{A_{g,\nu}|A_{g,x}})
=\operatorname{Norm}_{A_g}(A_{g,\nu}^!)=\operatorname{Norm}_{A_g}(A_{g,\nu}^{\circ}).
\]

Let $G$ be a topological group, equipped
with a conjugacy class \unsure{$\mathcal E$|$\mathcal X$} of subgroups,
\deleted{Let} $G_1$ in this class. We say that

\scanpage{209}
$G,\mathcal X$ is a Teichmüller \deleted{\unsure{datum}} \inserted{pair}
of type $g$, if there exists an isomorphism of
topological groups $G\simeq A_g$, taking
\unsure{$\mathcal E$|$\mathcal X$} to the conjugacy class of $A_1$.
\inserted{\unsure{It amounts to the same thing to say that}}
\deleted{We then see that $G_1=\operatorname{Norm}_G(G_1)$}
\deleted{hence that} $\mathcal X$, \deleted{\unsure{as}}
\inserted{equipped with the topology of a} homogeneous space
under $G$, is \deleted{homeomorphic to $G/G_1$} a
compact orientable surface of genus $g$,
\deleted{\illegible{2}} say \unsure{$\mathcal X$|$\mathfrak X$}, and that
\deleted{\illegible{3}} \inserted{the map}
natural
\[
(28)\qquad G\longrightarrow\operatorname{Aut}(\mathcal X)
\]
is a homeomorphism of topological groups.
We then see that \inserted{\unsure{by} $(G,\mathcal X)\longrightarrow\mathcal X$ from}
the category of
Teichmüller pairs of type $g$, to
that of compact oriented surfaces, is
an equivalence of categories. It follows
that \deleted{\unsure{every automorphism}} for an automorphism \inserted{\unsure{$u$}} of the
topological group $G$, $u$ is inner iff it

\scanpage{210}
\mspage{105}
preserves the class $\mathcal X$ i.e. iff $u(G_1)$ is conjugate
to $G_1$. Moreover the center of $G$ is trivial, \unsure{i.e.}

We can give an analogous description
for the category of oriented surfaces of
genus $g$ with $\nu$ holes (i.e. homeomorphic
to $U_{g,\nu}$) in terms of a topo-
logical group $G_\nu(\simeq A_{g,\nu})$ and of a \inserted{conjugacy class \unsure{$U_\nu$|$\widetilde U_\nu$} of}
subgroups $H_\nu$ of it $(\simeq A_{g,\nu}\cap A_{g,\nu+1})$, with the
conditions that $(G_\nu,\{H_\nu\})$ be isomorphic
to $(A_{g,\nu},A_{g,\nu}\cap A_{g,\nu+1}\unsure{)|\}}$, or else
that $U_\nu$ \inserted{(\unsure{space} \illegible{1} topology \unsure{of a homogeneous space})} be an oriented surface of
type $(g,\nu)$, and that the \deleted{map}canonical homomorphism
\[
(29)\qquad G\longrightarrow\operatorname{Aut}(U_\nu)
\]
be a homeomorphism of topological
groups. We then find an
equivalence between the category of
Teichmüller pairs $(G_\nu,\{H_\nu\})$
of type $(g,\nu)$, and that of
orientable surfaces of type $(g,\nu)$. We find again
